\chapter{Catálogos finitos y certificados de reproducibilidad}
\label{app:datos-reproducibles}

Los objetos finitos siguientes permiten repetir las verificaciones sin
incorporar notación de producción a la exposición matemática. Cada control
conserva el dominio, la operación, el invariante y el falsador que le
corresponden.

\section{Catálogos finitos del estado común}

\begin{longtable}{@{}>{\raggedright\arraybackslash}p{.34\textwidth}
  >{\raggedright\arraybackslash}p{.58\textwidth}@{}}
\toprule
Objeto finito & Contenido matemático verificable\\
\midrule
\endhead
Catálogo de emisiones de longitud seis
& Las \(468\) emisiones \(U_6\), sus multiplicidades, orientación y firmas
aditiva y multiplicativa.\\
Partición espectral de las emisiones
& Las veintiséis clases de tamaño \(18\), junto con el mapa explícito que
reconstruye cada emisión desde su clase y su posición interna.\\
Atlas dodecafásico de rutas
& Las \(468\) rutas mínimas con hoja, fase, acarreo, frontera y memoria.\\
Catálogo del lector arquimediano
& Las cinco regiones de \(\pi\), las regiones conjugadas de \(e\) y el eje de
autoescala de \(\varphi\), con sus cilindros y multiplicidades.\\
Flujo de Hensel de veinte bloques
& La matriz unimodular, su inversa entera, los estados \(3\)-ádicos y la
recurrencia de residuo--cociente que prolonga los tres canales.\\
Registro de monodromía nonádica
& Los veinte bloques, sus fases, prefijos ternarios y prefijos decimales;
incluye los \(43\) retornos completos y los \(387\) pares \((K,K+9)\) por
canal.\\
Ventana de supervivencia \(K=5,\ldots,13\)
& Los candidatos locales, la rama compatible y las tres salidas de la
transición \(9\to1\), sin identificar retorno de fase con reinicio del estado.\\
Lectura arquimediana inicial
& Las doce ternas decimales obtenidas de los primeros trece bloques de cada
canal, separadas del certificado de precisión arbitraria.\\
\bottomrule
\end{longtable}

\section{Identidades algebraicas, operatoriales y espectrales de los certificados finitos}

\begin{longtable}{@{}>{\raggedright\arraybackslash}p{.34\textwidth}
  >{\raggedright\arraybackslash}p{.58\textwidth}@{}}
\toprule
Control & Identidades y falsadores\\
\midrule
\endhead
Monodromía y generación de constantes
& Verifica \(w_6\to w_{30}\to R_{36}\to G_9\), la selección \(3\to1\),
el cambio de estado tras nueve fases, las cinco regiones de \(\pi\) y la
compatibilidad de truncamientos. Falla si se introduce una cifra objetivo
antes del lector.\\
Coordenada angular conjugada
& Comprueba el selector regional \(169\), la recurrencia determinantal, la
suma analítica y su cota de resto, y la unicidad de \(C^*\). Falla si una
constante del SI interviene en la selección.\\
Dinámica del cociente \(468=18\cdot26\)
& Certifica la biyección entre emisiones y clases, el retorno uniforme, la
medida de la sombra promediada, los cociclos de modo y la elevación
microscópica. Distingue la frecuencia cronológica de la medida de Parry.\\
Canales enteros de las constantes
& Verifica las identidades \((729,271,315,161)\), las realizaciones
conjuntistas de \(271/315\), las órbitas diedrales y la factorización de
paridad visible y paridad de acarreo.\\
Incidencia de Witt
& Reconstruye la incidencia punto--hexada, el módulo \(E_{30}\), el
transporte positivo de \(P_3\), la estrella de Witt y las obstrucciones de
equivarianza.\\
Selector variacional
& Reconstruye el normalizador \(S_3\), el promedio de Reynolds, la
factorización polar y la cadena de proyectores de rango tres; la positividad
se prueba exactamente en \(\mathbb Q(\sqrt5)\).\\
Elevación \(W_{12}\to W_{24}\)
& Comprueba las incidencias de cuatro y cinco puntos, la elevación única de
hexadas y la correspondencia tipada con las cinco regiones de \(\pi\).\\
Realizaciones de \(-1/12\)
& Construye los proyectores racionales de los desplazamientos de pasos tres
y cuatro, verifica simetría e idempotencia y obtiene
\((3-4)/12=-1/12\) como diferencia de trazas normalizadas.\\
Holonomía de Schläfli
& Recorre las plaquetas, certifica \(A^2=5I-4A+5J\), la configuración de las
veintisiete rectas, sus \(45\) triángulos y la partición
\(729=27+270+432\).\\
Funcionales espectrales
& Recalcula Catalán, Stieltjes, Apéry, Bendersky--Glaisher--Kinkelin, las
dos extensiones Euler--Kronecker y las cotas dirigidas de productos primos.\\
Vacancias y criticidad
& Comprueba las separaciones \(21/22\), los retornos \(6/7\), la parte
armónica finita y los aproximantes dodecafásicos del modo crítico.\\
\bottomrule
\end{longtable}

\section{Diccionario tipado de los catálogos}

\begin{longtable}{>{\ttfamily}p{.18\textwidth}>{\raggedright\arraybackslash}p{.72\textwidth}}
\toprule
Campo & Significado y codominio\\
\midrule
\endhead
Esig & Firma del canal multiplicativo; identifica una de las veintiséis
clases de la factorización del emisor.\\
hplus\_type & Clase de orientación del canal aditivo para los dos pares de
direcciones cardinales declarados.\\
diag\_c & Índice entero de la clase diagonal del cursor aditivo.\\
px\_size & Cardinal de la fibra de la clase aditiva correspondiente.\\
colw & Sexteto de pesos de columna del bloque emitido; pertenece a
\(\mathbb Z_{\geq0}^{6}\).\\
\bottomrule
\end{longtable}

\section{Reproducibilidad de los certificados y obstrucciones exactas de las construcciones finitas}

Los controles emplean aritmética entera o racional siempre que el objeto lo
permite. Las cotas analíticas usan intervalos con redondeo dirigido y separan
la demostración simbólica de la comparación decimal. Cada ejecución
reconstruye los objetos desde las semillas declaradas, verifica las identidades
de incidencia y naturalidad, y contrasta las salidas con una implementación
aritmética independiente.

El control del producto nativo rechaza cualquier dependencia de una evaluación
decimal, una factorización externa, una criba, una tabla de ceros o una
constante objetivo. El control nonádico rechaza la identificación de la fibra
ambiente de \(729\) elementos con el cardinal de un espacio de historias. El
control de fase rechaza \(g(K+9)=g(K)\) como prueba de igualdad de estados. El
control metrológico rechaza toda selección de ruta realizada después de
consultar una magnitud observada.

Una discrepancia entre una identidad del texto, el catálogo finito y su control
independiente refuta la certificación correspondiente. Esta regla no sustituye
ninguna demostración: fija un segundo canal de comprobación para los objetos
finitos que intervienen en ella.
